\chapter{Introduction to Artificial Intelligence}

\section{What is Artificial Intelligence}

The concept of Artificial Intelligence (AI) emerged shortly after World War II, with the term officially coined in 1956, and it continues to evolve today. AI encompasses a wide range of subfields, from general capabilities such as learning and perception to highly specific tasks like playing chess, proving mathematical theorems, generating poetry, driving in traffic, or diagnosing diseases~\parencite{RussellNorvig2010}. 

Definitions of AI vary, but they could be often categorized along two dimensions: human-centered (thinking and acting humanly) and rationality-centered (thinking and acting rationally). The human-centered perspective focuses on simulating human thought and behavior, exemplified by the Turing Test, which evaluates whether a machine can imitate human responses convincingly~\parencite{Turing1950}. The rationality-centered perspective emphasizes achieving optimal outcomes or correct reasoning regardless of human behavior, forming the basis for designing autonomous rational agents. Modern AI integrates these perspectives with data-driven learning to tackle complex real-world problems~\parencite{RussellNorvig2010}.

\section{Main Areas of AI}

Artificial Intelligence (AI) encompasses a wide range of subfields, each focused on enabling machines to perform tasks that typically require human intelligence. AI includes Machine Learning, Deep Learning, Computer Vision, Natural Language Processing, Interaction with Handwriting, Smart Robots and Interaction with Spoken Audio.

Machine Learning (ML) is a core area where software learns from data, identifies patterns, classifies information, makes predictions, and continuously improves through feedback mechanisms. A specialized branch of ML, Deep Learning, relies on artificial neural networks to handle complex tasks autonomously, processing data through multiple nonlinear layers to achieve high accuracy.

Other important areas include Computer Vision, which allows machines to interpret images and videos similarly to humans, and Natural Language Processing (NLP), which focuses on understanding and generating human language, supporting applications such as machine translation, sentiment analysis, and summarization. AI also addresses human-computer interaction, including handwriting recognition and spoken audio processing, enabling accurate interpretation of handwritten text or speech even under challenging conditions. Finally, intelligent robotics combines perception, learning, and adaptive behavior to perform a variety of tasks in dynamic environments~\parencite{Bendida2025}.

Together, these subfields illustrate the broad and versatile capabilities of AI technologies.